% Fase 6 — conclusões, limitações e trabalhos futuros.
% Base textual: monografia de Fernando Ferreira da Silva, cap. 7, pp. 131--133.
% Não se afirma validação experimental adicional às execuções descritas na dissertação.
\setlength{\parskip}{4pt}

\begin{frame}[plain]
\vspace*{1.84cm}
{\LARGE\color{navy}Síntese conclusiva\par}
\vspace{.25cm}
{\large Contribuições, limitações e perspectivas\par}
\vspace{.30cm}
\textcolor{navy!40}{\rule{.76\linewidth}{.6pt}}\par
\vspace{.22cm}
O que o MoeaBench organiza, o que os experimentos evidenciam e quais questões permanecem abertas.
\end{frame}

\begin{frame}{A contribuição principal: uma infraestrutura de avaliação}
\textbf{O trabalho não propõe um novo algoritmo evolutivo multiobjetivo.}
\par\medskip
A contribuição é integrar, em um fluxo experimental comum, decisões e evidências que muitas vezes ficam dispersas entre scripts e ferramentas.
\vspace{.22cm}
\begin{center}
\begin{tikzpicture}[>=Latex, every node/.style={font=\fontsize{8.5}{10.6}\selectfont}, node distance=.25cm]
\node[draw=navy!65,rounded corners=2pt,fill=pale,minimum width=2.25cm,minimum height=.95cm,align=center] (a) {Problema e\\algoritmo};
\node[draw=navy!65,rounded corners=2pt,fill=pale,minimum width=2.25cm,minimum height=.95cm,align=center,right=.34cm of a] (b) {Protocolo e\\execuções};
\node[draw=navy!65,rounded corners=2pt,fill=pale,minimum width=2.25cm,minimum height=.95cm,align=center,right=.34cm of b] (c) {Métricas e\\estatística};
\node[draw=navy!65,rounded corners=2pt,fill=pale,minimum width=2.25cm,minimum height=.95cm,align=center,right=.34cm of c] (d) {Relatórios e\\diagnósticos};
\draw[->,thick,navy!80] (a)--(b);
\draw[->,thick,navy!80] (b)--(c);
\draw[->,thick,navy!80] (c)--(d);
\end{tikzpicture}
\end{center}
\vspace{.15cm}
\begin{itemize}\setlength\itemsep{4pt}
\item O \textbf{experimento} preserva configurações, sementes, populações, histórico e metadados.
\item A API permite examinar resultados por \textbf{execução, geração, dominância e espaço} (decisões ou objetivos).
\item Persistência, visualização e relatórios viabilizam \textbf{reanálise} sem repetir a otimização.
\end{itemize}
\end{frame}

\begin{frame}{O que os exemplos efetivamente mostraram?}
\begin{columns}[T,onlytextwidth]
\begin{column}{.48\textwidth}
\textbf{Comparações entre MOEAs}\par\smallskip
NSGA-III e SPEA2 alcançaram valores finais próximos de HV em DTLZ2, apesar de trajetórias diferentes.
\par\bigskip
\textbf{Variabilidade estocástica}\par\smallskip
Cinco execuções produziram HV final estável, mas frentes com distribuições espaciais distintas.
\end{column}
\begin{column}{.48\textwidth}
\textbf{Normalização}\par\smallskip
Alterar a referência inverteu o ranking de HV bruto entre populações de 20 e 100 indivíduos.
\par\bigskip
\textbf{Estrutura FAIR}\par\smallskip
Em DTLZ1, \emph{closeness} $=0{,}883$ e \emph{headway} $=0{,}999$ coexistiram com quatro Q-scores nulos.
\end{column}
\end{columns}
\vspace{.26cm}
\textbf{Síntese:} bom desempenho em uma dimensão não assegura cobertura, organização espacial ou estabilidade em outras.
\end{frame}

\begin{frame}{Comparabilidade exige protocolo e referência explícitos}
\begin{columns}[T,onlytextwidth]
\begin{column}{.48\textwidth}
\textbf{O que preservar}\par\smallskip
\begin{itemize}\setlength\itemsep{4pt}
\item Problema, número de objetivos e orçamento de avaliações.
\item Configuração do MOEA e sementes independentes.
\item Execuções repetidas e dispersão dos indicadores.
\item Ground truth, normalização e calibração dos Q-scores.
\end{itemize}
\end{column}
\begin{column}{.49\textwidth}
\textbf{O que não concluir}\par\smallskip
\begin{itemize}\setlength\itemsep{4pt}
\item Um vencedor \emph{universal} entre algoritmos.
\item Equivalência entre métricas físicas e Q-scores.
\item Superioridade causal com orçamentos experimentais distintos.
\item Boa cobertura apenas porque as soluções são não dominadas.
\end{itemize}
\end{column}
\end{columns}
\vspace{.10cm}
\textbf{Papel do MoeaBench:} tornar as evidências rastreáveis e interpretáveis, \emph{não} substituir o julgamento científico do pesquisador.
\end{frame}

\begin{frame}{Limitação de engenharia: paralelismo e memória}
\textbf{O problema observado:} paralelizar múltiplas execuções por threads Python não trouxe ganho relevante em tarefas intensivas de CPU.
\vspace{.25cm}
\begin{columns}[T,onlytextwidth]
\begin{column}{.47\textwidth}
\textbf{Threads Python}\par\smallskip
No CPython convencional, o \emph{Global Interpreter Lock} (GIL) limita a execução simultânea de código Python intensivo de CPU. As tarefas se alternam, sem acelerar proporcionalmente o cálculo.
\end{column}
\begin{column}{.49\textwidth}
\textbf{Processos independentes}\par\smallskip
Contornam essa restrição, mas tendem a duplicar bibliotecas, estados e dados, elevando o consumo de memória. A máquina pode esgotar a RAM.
\end{column}
\end{columns}
\vspace{.32cm}
\textbf{Decisão atual:} priorizar execução serial previsível em vez de paralelismo com custo elevado de memória.
\par\medskip
\textbf{Perspectiva:} portar rotinas críticas para C ou extensões nativas capazes de liberar o GIL quando apropriado, examinando compilação, portabilidade e manutenção.
\end{frame}

\begin{frame}{Limitações e perguntas metodológicas em aberto}
\begin{columns}[T,onlytextwidth]
\begin{column}{.49\textwidth}
\textbf{Resolução e referência}\par\smallskip
\begin{itemize}\setlength\itemsep{5pt}
\item As frentes são amostradas por conjuntos \emph{finitos}.
\item Métricas de distância e calibrações dependem do GT e de sua distribuição.
\item Q-scores são relativos aos respectivos baselines, não uma garantia universal de qualidade.
\end{itemize}
\end{column}
\begin{column}{.49\textwidth}
\textbf{Generalização do diagnóstico}\par\smallskip
\begin{itemize}\setlength\itemsep{5pt}
\item Ampliar a validação para mais MOPs, MOEAs e cardinalidades.
\item Examinar frentes desconectadas ou degeneradas.
\item Quantificar a sensibilidade às amostragens e a incerteza dos escores.
\end{itemize}
\end{column}
\end{columns}
\vspace{.16cm}
\textbf{Consequência:} FAIR oferece evidências multidimensionais, mas a confiabilidade dessas evidências também precisa ser estudada experimentalmente.
\end{frame}

\begin{frame}{Perspectivas de continuidade}
\begin{columns}[T,onlytextwidth]
\begin{column}{.48\textwidth}
\textbf{Validação experimental}\par\smallskip
Ampliar baterias de MOEAs e MOPs, com diferentes objetivos, orçamentos e repetições.
\par\bigskip
\textbf{Calibração FAIR}\par\smallskip
Investigar a sensibilidade à densidade do GT, aos baselines e à geometria das frentes.
\end{column}
\begin{column}{.48\textwidth}
\textbf{Incerteza diagnóstica}\par\smallskip
Quantificar a confiança nos Q-scores e comunicar casos de evidência insuficiente.
\par\bigskip
\textbf{Engenharia}\par\smallskip
Estudar paralelismo por componentes nativos, sob restrições de memória e portabilidade.
\end{column}
\end{columns}
\vspace{.24cm}
\textbf{Prioridade científica:} testar a consistência das conclusões quando mudam a amostragem, o problema e o protocolo.
\end{frame}

\begin{frame}{Mensagem final}
\vspace{.31cm}
{\large\color{navy}\textbf{Avaliar um MOEA exige mais que executar um algoritmo e consultar uma métrica.}\par}
\vspace{.39cm}
O MoeaBench reúne \textbf{protocolo, execuções, histórico, referências, métricas, estatística, visualização e diagnósticos} em um fluxo rastreável.
\par\medskip
As métricas FAIR acrescentam uma leitura da \textbf{estrutura da frente resultante}: proximidade, cobertura, lacunas, regularidade, equilíbrio e progresso podem contar histórias diferentes sobre a mesma busca.
\vspace{.37cm}
\textcolor{navy!45}{\rule{\linewidth}{.6pt}}
\par\vspace{.19cm}
\textbf{Contribuição central:} aproximar implementação computacional e rigor metodológico para avaliações de MOEAs mais transparentes, reproduzíveis e cientificamente auditáveis.
\end{frame}
